\subsection{Polynomial Invariance and Block Separation}
\begin{proposition} \label{prop:kernel-range-polynomial-in-operator-invariant} \hlt{Kernel and Range of a Polynomial in an Operator Are Invariant}\\
Let $\mathcal{V}$ be a vector space over $\mathbb{F}$, let $T\in\LL(\mathcal{V})$, and let $p\in\mathbb{F}[z]$. Then:
\begin{enumerate}[label=(\roman*)]
\setlength{\itemsep}{0pt}
\item $\ker(p(T))$ is invariant under $T$;
\item $\operatorname{range}(p(T))$ is invariant under $T$.
\end{enumerate}
\begin{proof}
By Proposition \ref{prop:operator-polynomial-products}(ii), with $p$ and the polynomial $z$, the operators $p(T)$ and $T$ commute.
\begin{enumerate}[label=(\roman*)]
\setlength{\itemsep}{0pt}
\item Let $\mathbf{v}\in\ker(p(T))$, so $p(T)\mathbf{v}=\mathbf{0}$. By commutativity, $p(T)(T\mathbf{v})=T(p(T)\mathbf{v})=\mathbf{0}$; $T\mathbf{v}\in\ker(p(T))$, so $\ker(p(T))$ invariant under $T$.
\item Let $\mathbf{v}\in\operatorname{range}(p(T))$. By this membership, choose $\mathbf{u}\in\mathcal{V}$ such that $\mathbf{v}=p(T)\mathbf{u}$.\\
By commutativity, $T\mathbf{v}=Tp(T)\mathbf{u}=p(T)(T\mathbf{u})\in\operatorname{range}(p(T))$. Hence $\operatorname{range}(p(T))$ is invariant under $T$.
\end{enumerate}
\hlt{Structural Commentary}: Commutativity transfers $T$-stability separately to the polynomial kernel and range.\\
\hlt{Key Insight}: The identity $Tp(T)=p(T)T$ preserves both subspaces.
\end{proof}
\end{proposition}

\begin{theorem} \label{thm:sylvester-linear-matrix-equation} \hlt{Sylvester Linear Matrix Equation}\\
Let $m,n\geq1$, let $A\in\operatorname{Mat}(m,m,\mathbb{C})$, and let $B\in\operatorname{Mat}(n,n,\mathbb{C})$.\\
If $\operatorname{spec}(A)\cap\operatorname{spec}(B)=\emptyset$, then $(\forall C\in\operatorname{Mat}(m,n,\mathbb{C}))(\exists!X\in\operatorname{Mat}(m,n,\mathbb{C}))\,AX-XB=C$.
\begin{proof}
\begin{enumerate}[label=(\roman*)]
\setlength{\itemsep}{0pt}
\item \hlt{Sylvester Map}: Set $\mathcal{M}\defeq\operatorname{Mat}(m,n,\mathbb{C})$ and define $L:\mathcal{M}\to\mathcal{M}$ by $L(X)\defeq AX-XB$.\\
The map $L$ is linear by Proposition \ref{prop:matrix-product-associativity-and-identities}(iii)--(iv), and $\dim_{\mathbb{C}}(\mathcal{M})=mn$ by Proposition \ref{prop:matrix-spaces-are-vector-spaces}.\\
It suffices to prove $\ker(L)=\{0\}$ by Proposition \ref{prop:injective-iff-kernel-zero} and Corollary \ref{cor:equal-dimension-kernel-range}.
\item \hlt{Polynomial Intertwining}: Suppose $X\in\ker(L)$, so $AX=XB$. We prove $(\forall k\in\mathbb{Z}_{\geq0})\,A^kX=XB^k$. For $k=0$, this is $I_mX=XI_n$.\\
If $A^kX=XB^k$, then, by Proposition \ref{prop:matrix-product-associativity-and-identities} and $AX=XB$, $A^{k+1}X=A(A^kX)=A(XB^k)=(AX)B^k=(XB)B^k=XB^{k+1}$.\\
The claim follows by mathematical induction. Therefore $(\forall p\in\mathbb{C}[z])\,p(A)X=Xp(B)$ by Definition \ref{def:polynomial-evaluation-on-square-matrix}.
\item \hlt{Kernel Vanishing}: By Theorem \ref{thm:cayley-hamilton-matrix} and the polynomial-intertwining identity from (ii), $p_B(A)X=Xp_B(B)=0$.\\
If $\lambda\in\operatorname{spec}(A)$, then $\lambda\notin\operatorname{spec}(B)$ by the disjointness hypothesis, so $p_B(\lambda)\neq0$ by Theorem \ref{thm:matrix-characteristic-polynomial-basic-coefficients}(ii).\\
Thus $p_B(A)$ is invertible by Corollary \ref{cor:polynomial-in-matrix-invertible-spectrum-test}, and $X=p_B(A)^{-1}p_B(A)X=0$. Hence $\ker(L)=\{0\}$.
\item \hlt{Existence and Uniqueness}: The map $L$ is injective by (iii) and Proposition \ref{prop:injective-iff-kernel-zero}.\\
The map $L$ is surjective by Corollary \ref{cor:equal-dimension-kernel-range} and invertible by Proposition \ref{prop:linear-map-invertible-iff-bijective}.\\
Consequently, $(\forall C\in\mathcal{M})(\exists!X\in\mathcal{M})\,L(X)=C$, equivalently, $(\forall C\in\mathcal{M})(\exists!X\in\mathcal{M})\,AX-XB=C$.
\end{enumerate}
\hlt{Structural Commentary}: Polynomial intertwining combines Cayley--Hamilton for $B$ with spectral invertibility for $p_B(A)$.\\
\hlt{Key Insight}: Disjoint spectra force every homogeneous Sylvester solution to vanish.
\end{proof}
\end{theorem}

Figure \ref{fig:sylvester-equation-spectral-gap} summarizes the spectral-separation argument behind invertibility of the Sylvester map.

\begin{illustrationfigure}[H]
  \centering
  \begin{tikzpicture}[
  >=Stealth,
  every node/.style={font=\small},
  box/.style={draw=gray!55, fill=gray!5, rounded corners=2pt, align=center, inner xsep=5pt, inner ysep=4.5pt},
  good/.style={draw=blue!60!black, fill=blue!7, rounded corners=2pt, align=center, inner xsep=5pt, inner ysep=4.5pt},
  arr/.style={->, line width=0.75pt, blue!70!black}
]
  \node[box] (gap) at (0,0) {$\operatorname{spec}(A)\cap\operatorname{spec}(B)=\emptyset$};
  \node[good] (poly) at (3.61,0) {$p_B(A)$ invertible};
  \node[box] (map) at (6.78,0) {$L^{-1}$ exists};

  \draw[arr] (gap.east) -- (poly.west);
  \draw[arr] (poly.east) -- (map.west);
\end{tikzpicture}

  \caption{Disjoint spectra make $p_B(A)$ invertible, so the Sylvester map $L(X)=AX-XB$ is invertible.}
  \label{fig:sylvester-equation-spectral-gap}
\end{illustrationfigure}

\begin{corollary} \label{cor:block-upper-triangular-remove-disjoint-coupling} \hlt{Removing a Spectrally Separated Off-Diagonal Block}\\
Let $m,n\geq1$, let $A\in\operatorname{Mat}(m,m,\mathbb{C})$, let $B\in\operatorname{Mat}(n,n,\mathbb{C})$, and let $C\in\operatorname{Mat}(m,n,\mathbb{C})$. Set $M\defeq\left(\begin{smallmatrix}A&C\\0&B\end{smallmatrix}\right)\in\operatorname{Mat}(m+n,m+n,\mathbb{C})$.\\
If $\operatorname{spec}(A)\cap\operatorname{spec}(B)=\emptyset$, then $M$ is similar to $A\oplus B$.
\begin{proof}
By Theorem \ref{thm:sylvester-linear-matrix-equation}, choose $X\in\operatorname{Mat}(m,n,\mathbb{C})$ such that $AX-XB=-C$. Set $P\defeq\left(\begin{smallmatrix}I_m&-X\\0&I_n\end{smallmatrix}\right)$ and $Q\defeq\left(\begin{smallmatrix}I_m&X\\0&I_n\end{smallmatrix}\right)$.\\
By Proposition \ref{prop:block-matrix-arithmetic}(ii), $PQ=QP=I_{m+n}$, so $P$ is invertible and $P^{-1}=Q$. By same proposition, $PMQ=\left(\begin{smallmatrix}A&AX+C-XB\\0&B\end{smallmatrix}\right)$.\\
By $AX-XB=-C$, $PMQ=\left(\begin{smallmatrix}A&0\\0&B\end{smallmatrix}\right)=A\oplus B$. Thus $PMP^{-1}=A\oplus B$ by $P^{-1}=Q$, and $M$ is similar to $A\oplus B$ by Definition \ref{def:similar-matrices}.\\
\hlt{Structural Commentary}: A Sylvester solution determines the block shear that removes the coupling.\\
\hlt{Key Insight}: Spectral separation makes the upper-right block removable by similarity.
\end{proof}
\end{corollary}

\begin{theorem} \label{thm:schur-block-form-by-eigenvalue-groups} \hlt{Schur Form Grouped by Eigenvalues}\\
Let $n,d,n_1,\ldots,n_d\geq1$, let $\lambda_1,\ldots,\lambda_d\in\mathbb{C}$ be distinct, and let $A\in\operatorname{Mat}(n,n,\mathbb{C})$.\\
Suppose $\sum\limits_{i=1}^{d}n_i=n$, $\operatorname{spec}(A)=\{\lambda_1,\ldots,\lambda_d\}$, and $(\forall i\in\{1,\ldots,d\})\,m_A(\lambda_i)=n_i$.
\begin{enumerate}[label=(\roman*)]
\setlength{\itemsep}{0pt}
\item The matrix $A$ is unitarily similar to a block upper-triangular matrix
\begin{align}
T=
\begin{bmatrix}
T_{11}&T_{12}&\cdots&T_{1d}\\
0&T_{22}&\cdots&T_{2d}\\
\vdots&\vdots&\ddots&\vdots\\
0&0&\cdots&T_{dd}
\end{bmatrix},
\qquad
T_{ii}\in\operatorname{Mat}(n_i,n_i,\mathbb{C}),
\nonumber
\end{align}
where each $T_{ii}$ is upper triangular with every diagonal entry equal to $\lambda_i$.
\item The matrix $A$ is similar to $T_{11}\oplus\cdots\oplus T_{dd}$.
\end{enumerate}
\begin{proof}
\begin{enumerate}[label=(\roman*)]
\setlength{\itemsep}{0pt}
\item Order the eigenvalue multiset of $A$ by listing $\lambda_i$ exactly $n_i$ times, $i\in\{1,\ldots,d\}$. Choose a unit vector $\mathbf{x}\in E_{\lambda_1}(A)$.\\
By Theorem \ref{thm:matrix-schur-triangularization-ordered}(i), $T\defeq U^*AU$ is upper triangular with this ordered diagonal for some unitary $U$.\\
After grouping consecutive rows and columns according to $(n_1,\ldots,n_d)$, $T$ has the stated block form and each $T_{ii}$ has constant diagonal $\lambda_i$. Thus $A$ is unitarily similar to $T$ by Definition \ref{def:unitary-and-permutation-similarity}.
\item By part (i), Theorem \ref{thm:unitary-matrix-equivalent-conditions}(vi), and Definition \ref{def:similar-matrices}, $A$ is similar to $T$. If $d=1$, then $T=T_{11}$.\\
Suppose $d\geq2$. Set $R_0\defeq T$, and for $k\in\{1,\ldots,d-1\}$ let $R_k$ be the trailing principal block with block indices $k+1,\ldots,d$. Then $R_{k-1}=\left(\begin{smallmatrix}T_{kk}&C_k\\0&R_k\end{smallmatrix}\right)$. By Proposition \ref{prop:upper-triangular-characteristic-polynomial}, $\operatorname{spec}(T_{kk})=\{\lambda_k\}$ and $\operatorname{spec}(R_k)=\{\lambda_{k+1},\ldots,\lambda_d\}$.\\
Hence $R_{k-1}$ is similar to $T_{kk}\oplus R_k$ by Corollary \ref{cor:block-upper-triangular-remove-disjoint-coupling}.\\
Consequently, for every $k\in\{1,\ldots,d-1\}$, $T$ is similar to $T_{11}\oplus\cdots\oplus T_{kk}\oplus R_k$ by mathematical induction, Proposition \ref{prop:direct-sum-matrix-algebra}(ii), and transitivity in Proposition \ref{prop:similarity-equivalence-relation}. Since $R_{d-1}=T_{dd}$, the conclusion follows by transitivity in Proposition \ref{prop:similarity-equivalence-relation}.
\end{enumerate}
\hlt{Structural Commentary}: Ordered Schur form groups eigenvalues; Sylvester shears then split successive spectral blocks.\\
\hlt{Key Insight}: Cross-group blocks are coordinate-dependent and removable by similarity.
\end{proof}
\end{theorem}

\begin{corollary} \label{cor:commuting-with-disjoint-blocks-has-zero-off-diagonal-blocks} \hlt{Commuting with Disjoint Blocks}\\
Let $d,n_1,\ldots,n_d\geq1$ and let $A_{ii}\in\operatorname{Mat}(n_i,n_i,\mathbb{C})$, $i\in\{1,\ldots,d\}$, have pairwise disjoint spectra.\\
Set $n\defeq\sum\limits_{i=1}^{d}n_i$ and $A\defeq A_{11}\oplus\cdots\oplus A_{dd}\in\operatorname{Mat}(n,n,\mathbb{C})$.\\
If a conformally partitioned $B=[B_{ij}]\in\operatorname{Mat}(n,n,\mathbb{C})$ satisfies $AB=BA$, then $B_{ij}=0$ for all distinct $i,j\in\{1,\ldots,d\}$.
\begin{proof}
By $AB=BA$ and Proposition \ref{prop:block-matrix-arithmetic}(ii), $(\forall i,j\in\{1,\ldots,d\})\,A_{ii}B_{ij}-B_{ij}A_{jj}=0$.\\
Fix $i,j\in\{1,\ldots,d\}$ with $i\neq j$. The spectra of $A_{ii}$ and $A_{jj}$ are disjoint.\\
Hence $B_{ij}=0$ by the uniqueness assertion in Theorem \ref{thm:sylvester-linear-matrix-equation} with zero right-hand side.\\
\hlt{Structural Commentary}: Block commutation reduces every off-diagonal entry to a homogeneous Sylvester equation.\\
\hlt{Key Insight}: Homogeneous spectral separation forces each off-diagonal block to vanish.
\end{proof}
\end{corollary}
